\begin{claim}{D.53}{Proved}{The unique labelled order-four section class}
Let $a$ have order four in the full punctured-torus
loop-outer group, and let
\[
          1\longrightarrow F\longrightarrow E
             \xrightarrow{\;\pi\;}C_4=\langle a\rangle
             \longrightarrow1,\qquad F=\widehat F_2,
\]
be its centre-free extension. Its homomorphic
sections form exactly one $F$-conjugacy class,
with the generator label $a$ retained.
For an actual lift $b\in E$ with
$b^4=1$ and $\pi(b)=a$, put $B=\langle b\rangle$.
Then
\[
       N_E(B)=C_F(b)\times B,\qquad
                         H^1(C_F(b),\F_2)=0.            \tag{D.53a}
\]
Conjugation by a lift of $a$ fixes exactly one
of the three full section classes for $a^2$
and interchanges the other two.
Consequently $E$ has two conjugacy classes
of subgroups of order two; exactly one of
them is contained in its class of cyclic
subgroups of order four.
\end{claim}
\begin{proof}
D.46 supplies the actual lift $b$ with its
prescribed label. Put $T=\langle b^2\rangle$
and
\[
               E_2=\pi^{-1}(\langle a^2\rangle)
                      =F\rtimes T.
                                                               \tag{D.53b}
\]
The action of $b^2$ on $H_1(F,\Z_2)$ is
$-I$ by D.46. Thus
\[
                      E_2^{\rm ab,(2)}\cong C_2^3.
                                                               \tag{D.53c}
\]
Use coordinates $(v,u)\in\F_2^2\times\F_2$,
where $u$ is the quotient character and
the first two coordinates are the extended
characters $\chi_1,\chi_2$ of D.46.
The three full section classes of $E_2$
give three distinct points of the affine
plane $u=1$.
Indeed, the degree-zero localized algebra
in D.46g--i has dimension three, is generated
by $\chi_1/u,\chi_2/u$, and restriction to
the three section normalizers identifies
it with $\F_2^3$.
Each restriction is evaluation at
$(\chi_1(t),\chi_2(t))$ for the nonidentity
element $t$ of that section.
If two such points agreed, the two
restriction coordinates would agree on
both generators and hence on the entire
algebra, contradicting the stated product
isomorphism.

Let $R=A_2(a)\in\GL_2(\Z_2)$.
D.50 gives $R^2=-I$.
Moreover $R\not\equiv I\pmod2$.
For otherwise $R=I+2D$ with an integral
matrix $D$, so $R^2=-I$ would give
$I+2D+2D^2=0$, an impossibility modulo two.
Conjugation by $b$ on (D.53c) fixes $u$
and induces the actual affine map
\[
               v\longmapsto\overline R v+d
                        \quad\text{on }u=1             \tag{D.53d}
\]
for a vector $d\in\F_2^2$.
The linear part follows from the
restriction of conjugation to the
specified subgroup $F$; the translation
is the image of a chosen point of $u=1$.
Conjugation by an element of $E_2$
acts trivially on the three
$F$-conjugacy classes, as shown in D.46.
Thus (D.53d) permutes their three points
with order dividing two.
Three distinct points of $\F_2^2$ affinely
span it. Fixing all three would therefore
make this affine map the identity,
contrary to $\overline R\ne I$.
Its permutation is consequently a
transposition, with exactly one fixed class.
The class of $T=\langle b^2\rangle$ is
that fixed class.

Put $C=C_F(T)$ and $N=N_E(T)$.
Since $b$ normalizes $T$, the image of $N$
under $\pi$ is all of $C_4$.
Its kernel is exactly $C$, because an
element of $F$ normalizing $T$ centralizes
its nonidentity element.
The chosen section gives the actual
semidirect-product identification
\[
          N=C\rtimes\langle b\rangle,\qquad
          T\subseteq Z(N),\qquad
          J:=N/T=C\rtimes\langle\overline b\rangle,
                    \quad|\overline b|=2.              \tag{D.53e}
\]
In the last expression the embedding of
$C$ is injective since $C\cap T=1$.
Conjugation by $b$ on $C$ has square the
identity because $b^2$ centralizes $C$.

As a closed subgroup of the full free
profinite group $F$, the group $C$ is
projective and has $\operatorname{cd}_2(C)\leq1$.
D.46 gives $H^1(C,\F_2)=0$.
All its positive-degree mod-two cohomology
therefore vanishes.
Hochschild--Serre for the last extension
in (D.53e) identifies inflation as
\[
               H^*(C_2,\F_2)\xrightarrow{\;\cong\;}
                                  H^*(J,\F_2).         \tag{D.53f}
\]
This uses the displayed full group $C$,
not an identification of $C$ with its
maximal pro-two quotient.

We spell out the full Brown hypotheses
for this new normalizer quotient.
The subgroup $C$ is open in $J$ and has
finite $2$-cohomological dimension.
It is torsion-free. Every finite subgroup
of $J$ consequently injects into $J/C=C_2$;
in particular $J$ has no subgroup $C_2^2$.
Theorem~2 of
\cite[\S1.1]{A-KochloukovaZalesskii2007}
therefore supplies, in degrees $n>1$,
the dense restriction embedding into the
product over the conjugacy classes of
order-two subgroups of $J$.
By (D.53f) its source has dimension one.
Each product factor is nonzero: for a
section $S$ of $J\to C_2$, its normalizer is
\[
                        C_C(S)\times S,
                                                               \tag{D.53g}
\]
and the restriction from the $S$ factor
already supplies one dimension.
Thus the product has finitely many
factors and the finite dense image is
the entire product.
Since $\operatorname{cd}_2(C_C(S))\leq1$,
the dimensions give the exact equality
\[
             1=\sum_{[S]}
                       \bigl(1+\dim H^1(C_C(S),\F_2)\bigr).
                                                               \tag{D.53h}
\]
There is a section $\langle\overline b\rangle$.
It follows that it is the unique
conjugacy class and that
$H^1(C_C(\overline b),\F_2)=0$.
These are also $C$-conjugacy classes:
an element of $J$ can be multiplied by
the appropriate element of a given
section to leave a conjugator in $C$.

Now let $B'$ be any subgroup of order
four in $E$. Its intersection with $F$
is trivial, so it is cyclic and $\pi|_{B'}$
is an isomorphism onto $C_4$.
Choose its unique generator $b'$ with
$\pi(b')=a$. Its square generates an
order-two section $T'$ of $E_2$.
Conjugation by $b'$ has the same action
on the section classes of $E_2$ as
conjugation by $b$, because
$b'b^{-1}\in F$.
The class of $T'$ is fixed, and is
therefore the unique class of $T$.
An $F$-conjugation puts $T'$ equal to $T$,
and then $B'\subseteq N_E(T)=N$.
Its image in $J=N/T$ has order two.
By (D.53h), a further conjugation by
an element of $C$ puts that image equal
to $\langle\overline b\rangle$.

The exact inverse image of this subgroup
under $N\to N/T$ has order four,
because its kernel is the specified
subgroup $T$ of order two.
It contains $\langle b\rangle$ and
hence equals $\langle b\rangle$.
Thus the conjugated $B'$ equals $B$.
All conjugations used lie in $F$ and
preserve $\pi$, so its labelled generator
$b'$ is sent to $b$, proving uniqueness
of the homomorphic section with labels.

The fixed group $C_C(\overline b)$ in
(D.53h) is exactly $C_F(b)$.
For an element of $C$ fixed by
$\overline b$, its commutator with $b$
lies in $T$ and also in $C$, and
$T\cap C=1$.
The converse is immediate, and an
element commuting with $b$ already
centralizes $T$.
This proves the cohomology statement
in (D.53a).
Finally, an element normalizing $B$
must fix its generator $b$: conjugation
preserves its label in the abelian
quotient $C_4$, and $\pi|_B$ is injective.
Multiplication by the matching power
of $b$ then leaves an element of
$C_F(b)$, proving the direct-product
normalizer in (D.53a).

The two orbits of the transposition
in (D.53d) are exactly the two
$E$-conjugacy classes of order-two
subgroups. Only its fixed class can
be the square of a subgroup of order
four, and $B$ realizes that containment.
This proves the final assertion.

The finite program
\path{code/part_c/torus_order4_sections_d53.py}
checks the mod-four matrix obstruction
and every invariant three-point affine
configuration above a nonidentity
order-two matrix in $\GL_2(\F_2)$.
The normalizer and restriction arguments
prove the full section result.
No conjugacy of the outer action $a$
to a prescribed geometric action is
asserted here.
\end{proof}
